% !TeX root = ../../thesis.tex
\chapter{Virtual load and strain sensing overview}\label{ch:Virtual_load_and_strain_sensing_overview}

Start with general equations and concepts on dynamic system models, state estimation, and joint input-state estimation. Silvia Vettori's text is useful here.

\section{Dynamic system models}

Dynamic system models are mathematical representations of physical systems that evolve over time. They are used to describe the behavior of systems in various fields, including engineering, physics, and economics. These models can be classified into different categories, such as linear and nonlinear models, discrete and continuous models, and deterministic and stochastic models.

\begin{enumerate}
    \item state space models
    \item system models (maybe separating the second line of the state space representation?)
\end{enumerate}

\section{State estimation}

\begin{enumerate}
    \item Should have something about probability
    \item Kalman filters
\end{enumerate}

\section{Joint input-state estimation}

State estimation is the process of estimating the internal state of a dynamic system based on available measurements. In many cases, the system's state is not directly observable, and the goal is to infer it from noisy and incomplete measurements. Joint input-state estimation extends this concept by simultaneously estimating the system's state and the inputs acting on it.



\ldots

\begin{enumerate}
    \item this chapter is for estimation!
\end{enumerate}


%%%%%%%%%%%%%%%%%%%%%%%%%%%%%%%%%%%%%%%%%%%%%%%%%%
% Keep the following \cleardoublepage at the end of this file, 
% otherwise \includeonly includes empty pages.
\cleardoublepage

